\section{Errors that are data}
\label{sec:ch05-errors-that-are-data}
\index{mutations!conventions|(}

Four chapters of a read-only graph, and here is the first write. It arrives now
rather than in chapter~\ref{ch:hotchocolate-quickly}, where it would have been
easy, because a mutation with nothing to refuse teaches nothing.

\begin{minted}[fontsize=\footnotesize]{csharp}
[Error<ProductNotFoundError>]
[Error<CustomerNotFoundError>]
[Error<RatingOutOfRangeError>]
[Error<DuplicateReviewError>]
public static async Task<Review> SubmitReviewAsync(
    [ID<Product>] Guid productId,
    [ID<Customer>] Guid customerId,
    int rating,
    string? body,
    CatalogService catalog,
    AccountsService accounts,
    ReviewsService reviews,
    ITopicEventSender sender,
    CancellationToken cancellationToken)
\end{minted}

A method with loose parameters returning a \code{Review}. Here is what the
schema makes of it:

\begin{graphqlsdl}
type Mutation {
  submitReview(input: SubmitReviewInput!): SubmitReviewPayload!
}

type SubmitReviewPayload {
  review: Review
  errors: [SubmitReviewError!]
}
\end{graphqlsdl}

\index{mutations!AddMutationConventions}%
One call did that reshaping, registered on the builder with its
\code{applyToAllMutations} flag set:

\begin{minted}{csharp}
.AddMutationConventions(applyToAllMutations: true)
\end{minted}

What it produces is the shape every mutation in this book will have from here:
one argument called \code{input}, one payload, and the payload carrying the
errors~\autocite{chillicream2026mutations}. Services and the
\code{CancellationToken} are left out of the generated input, because only
parameters that map to GraphQL arguments belong there.

The \code{applyToAllMutations} flag is the decision worth defending. The
alternative is an attribute on each mutation, applied by whoever remembers.
Turning it on globally means the schema cannot drift into two shapes depending
on who wrote which field, and the escape hatch for the case that genuinely
differs is then explicit rather than accidental.

\subsection*{Why the payload has an input type at all}

\index{mutations!input types}%
The single \code{input} argument looks like ceremony until you have needed it.
Arguments on a field are positional in the sense that matters: adding one is
additive, but you cannot deprecate a required one, you cannot reorder them
meaningfully, and a client library generating a method signature has to
regenerate every time the list changes. An input object is a type, so it evolves
under the same rules every other type does, including the ones
section~\ref{sec:ch05-what-a-client-can-survive} spent its length on.

The payload is the same argument on the way out. A mutation that returns
\code{Review} directly has nowhere to put anything else, ever. A mutation that
returns \code{SubmitReviewPayload} can grow a field later without touching the
one thing it already returns, and this one already has: \code{errors} is a field
that exists only because the payload does.

\subsection*{Four ways to be told no}

\index{errors!typed domain errors}%
The four attributes at the top of the method are the whole error model, and this
is what they generate:

\begin{graphqlsdl}
union SubmitReviewError =
  | ProductNotFoundError
  | CustomerNotFoundError
  | RatingOutOfRangeError
  | DuplicateReviewError
\end{graphqlsdl}

A union is Mosaic's second abstract type and the counterpart to the interface in
section~\ref{sec:ch05-one-name-for-one-thing}. An interface says these types
share a field; a union says a value is one of these types and makes no claim
about what they have in common. Which is the right relationship here: the four
errors have nothing in common as data, and what they share is that exactly one
of them can come back.

\index{HotChocolate!Error attribute}%
Two notes about the attribute. \code{[Error<T>]} is in the source but not in the
documentation, which shows only \code{[Error(typeof(T))]}. They do the same
thing, and Mosaic uses the generic form because the rest of the codebase already
reads \code{[ObjectType<T>]} and \code{[QueryType]}.

More important is what \code{T} points at. It can be the exception itself, which
is what the documentation shows, and then the exception's \code{Message} becomes
the payload's message~\autocite{chillicream2026mutations}. Mosaic points it at
an error class instead:

\begin{minted}[fontsize=\footnotesize]{csharp}
public sealed class RatingOutOfRangeError : IMosaicError
{
    private RatingOutOfRangeError(string message) => Message = message;

    public string Message { get; }

    public string Code => "RATING_OUT_OF_RANGE";

    public int Min => 1;

    public int Max => 5;

    public static RatingOutOfRangeError CreateErrorFrom(
        RatingOutOfRangeException exception)
        => new($"A rating has to be between 1 and 5. Got {exception.Rating}.");
}
\end{minted}

The static \code{CreateErrorFrom} is what the extra class buys. Pointing the
attribute at the exception copies a message written for a log into a field a
customer will read, and leaves nowhere at all to put a code. Here the exception
is an internal signal and the error is a public type, and the factory is the
seam between them.

\index{errors!fields beyond the message}%
\code{Min} and \code{Max} are the part I would push hardest. A client rendering
a form can read them and correct the input without knowing what Mosaic considers
a valid rating. That is not possible in the \code{errors} array at any price:
the array's entries have a fixed shape, and anything extra goes into
\code{extensions} as a bag with no type. Here it is two integers on a type, and
a client that generates code from the schema gets them as two integers.

It also answers a question the type system raises and cannot settle.
\code{rating} is an \code{Int!}, and GraphQL has no way to say that an integer
lies between one and five. Every constraint the type system cannot express has
to live somewhere, and the only choice is whether it lives somewhere the client
can read.

\subsection*{An interface with a code on it}

\index{errors!custom interface}%
HotChocolate's default error interface carries one field, \code{message}.
Mosaic replaces it:

\begin{minted}[fontsize=\footnotesize]{csharp}
[GraphQLName("Error")]
public interface IMosaicError
{
    string Message { get; }

    string Code { get; }
}
\end{minted}

\begin{minted}{csharp}
builder.AddGraphQL()
    .AddErrorInterfaceType<IMosaicError>()
\end{minted}

\begin{graphqlsdl}
interface Error {
  message: String!
  code: String!
}

type RatingOutOfRangeError implements Error {
  message: String!
  code: String!
  min: Int!
  max: Int!
}
\end{graphqlsdl}

Error classes are not required to implement the C\# interface; matching
properties are enough~\autocite{chillicream2026mutations}. Mosaic implements it
anyway, so that a new error type which forgets its code fails to compile rather
than failing to build the schema.

\index{errors!codes as strings}%
The code is a \code{String} and not an enum, which is the less obvious choice
and I will defend it. An enum would be checked at build time and would let a
client switch exhaustively over the codes, which sounds better until you add the
fifth error. Every exhaustive client then has a case it has never seen, and
adding a domain error becomes a breaking change for the people most carefully
handling your errors. A string means an unrecognised code falls back to
displaying the message, which is what a client should do anyway with a failure
it does not understand.

\subsection*{What it looks like on the wire}

\index{errors!union selection}%
The union has a cost and it shows up the first time you write the query. This
does not work:

\begin{minted}[fontsize=\footnotesize]{text}
A union type cannot declare a field directly.
Use inline fragments or fragments instead.
\end{minted}

A union makes no claim that its members share anything, so nothing can be
selected on it directly, not even \code{message}. The working form leans on the
interface:

\begin{minted}[fontsize=\footnotesize]{graphql}
errors {
  __typename
  ... on Error { message code }
  ... on RatingOutOfRangeError { min max }
}
\end{minted}

One fragment covers every member, because every member implements \code{Error},
and a second fragment picks up the fields only one of them has. That first
fragment is the return on the custom interface: with the default one a client
would get \code{message} this way and would have no route to a code at all.

Four requests, all HTTP 200, all with \code{data} non-null:

\begin{minted}[fontsize=\footnotesize]{json}
{"data":{"submitReview":{
  "review":{"id":"UmV2aWV3Og3knwEW8Fx4r17UOtwmS5o=","rating":4,
            "body":"Solid, but the finish scratches.",
            "author":{"displayName":"Sofia Ferrari"}},
  "errors":null}}}
\end{minted}

\begin{minted}[fontsize=\footnotesize]{json}
{"data":{"submitReview":{"review":null,"errors":[
  {"__typename":"DuplicateReviewError",
   "message":"You have already reviewed this product.",
   "code":"DUPLICATE_REVIEW"}]}}}
\end{minted}

\begin{minted}[fontsize=\footnotesize]{json}
{"data":{"submitReview":{"review":null,"errors":[
  {"__typename":"RatingOutOfRangeError",
   "message":"A rating has to be between 1 and 5. Got 9.",
   "code":"RATING_OUT_OF_RANGE","min":1,"max":5}]}}}
\end{minted}

No \code{errors} key at the top level in any of them. A refused mutation is a
successful request that returned a refusal, and the client reads the refusal the
same way it reads the review: by selecting fields.

\index{errors!what stays in the array}%
Anything not named in one of the four attributes still goes to the array as an
unexpected execution error, and that is the arrangement working rather than a
gap in it. A bug is not a domain error. It has no branch behind it, the client
can do nothing with it, and the last thing it should do is arrive looking like a
considered answer.

\subsection*{A check with an expiry date}

\index{federation!referential integrity}%
Two of the four errors are on borrowed time and the code says so in a comment.

\begin{minted}[fontsize=\footnotesize]{csharp}
if (await catalog.GetProductByIdAsync(productId, cancellationToken) is null)
{
    throw new ProductNotFoundException(productId);
}
\end{minted}

Reviews can ask Catalog whether a product exists because Catalog is a folder in
the same process, holding tables in the same database, inside the same
transaction boundary. From
chapter~\ref{ch:the-first-cut-extracting-catalog} onward it is neither, and this
check becomes a network call on the write path or an accepted risk.

I am flagging it now rather than discovering it then, because the flinch this
produces is the thing worth having early. A referential check across a boundary
that does not exist yet is the most common way a monolith turns out not to have
been ready. \code{RatingOutOfRangeError} and \code{DuplicateReviewError} have no
such problem: Reviews owns both rules and both rows, and both survive the move
untouched. Notice that the two that survive are the two where the domain owns
the data the rule is about. That is not a coincidence and it is a good test to
apply to a rule before writing it.

\subsection*{Verify it in Postman}

\index{Postman!mutation assertions}%
The collection grew by nine requests this chapter, six of them for the material
in the sections above: refetching through \code{node}, a page of reviews, the
next page from its cursor, a product nobody has reviewed, the deprecated field
still answering, and the deprecation visible in introspection.

The last three are the mutation, and they run in this order for a reason. First
a rating of nine, which asserts the code, the message and both bounds without
writing anything. Then a review from a customer who has not reviewed that
product, which asserts the payload carries no errors, that the review came back
holding what was sent, and that its identifier decodes to a \code{Review} node.
Then the identical request again, which asserts \code{review} is null and that
the single error is \code{DuplicateReviewError} with code
\code{DUPLICATE\_REVIEW}.

\index{Postman!ordering}%
Order is load-bearing here for the first time in this book, and it changed how
the whole gate works. The mutation writes a row. Every assertion about the
seeded catalog, twenty-five products and a hundred and twenty reviews, has to
run before it, or the second run of the collection finds a hundred and
twenty-one and fails an assertion that is not wrong.

That was not enough on its own. A verification run now leaves the database
changed, so both scripts start the service with
\code{MOSAIC\_RESET\_DATABASE=1}, and the seeder drops the schema and reseeds
before the service takes a request. A gate whose result depends on how many
times it has been run is not a gate. The cost is a second or two and the check
that it worked is running \code{verify.ps1} twice in a row.

\begin{minted}[fontsize=\footnotesize]{text}
requests      19
assertions    74
\end{minted}

\index{diagnostics!configuration binding}%
One trap for anybody adding a switch like that. The obvious reading is the one
that binds the value straight to a boolean:

\begin{minted}[fontsize=\footnotesize]{csharp}
if (configuration.GetValue<bool>("MOSAIC_RESET_DATABASE"))
\end{minted}

That accepts only \code{true} and \code{false}. Given \code{1} it throws
\code{InvalidOperationException}, unhandled, and the host fails to start. Not a
fallback to false: the service does not come up. Mosaic's other switch is
\code{MOSAIC\_KEEP\_DATABASE=1}, so \code{=1} was the shape I copied, and the
seeder now reads the raw string and compares it by hand.

So the schema can refuse now, and say why, in terms a client can act on. What it
cannot do is start a sentence.
\index{mutations!conventions|)}
